\chapter{Experimental search for dark matter}
\label{ch:dm-search}

Direct laboratory searches for dark matter (DM) can be organised around
the type of particle one is looking for and the way it interacts in
the detector. In this chapter we treat two channels separately:
\emph{recoil}, in which a fermion-like DM particle scatters elastically
off a nucleus or a bound electron and deposits a small fraction of
its kinetic energy; and \emph{absorption}, in which a coherent
bosonic DM field is absorbed by the detector and deposits its full
rest energy at once. The two channels demand different kinematics
and different rate formulas, but they share the same downstream
machinery: the predicted spectrum has to be folded with detector
threshold, efficiency and resolution to produce a count, and the
count then has to be turned into a statistical statement about the
underlying model parameters. The chapter walks both channels through
the same four-stage pipeline --- \emph{model parameters} $\to$
\emph{predicted in-detector spectrum} $\to$ \emph{measured spectrum}
$\to$ \emph{statistical limit on the model parameter} --- producing
the formal apparatus that the four mechanism-specific chapters
that follow
(\cref{ch:phonon-qp,ch:ir-junction,ch:phonon-qubit-piezo,ch:microwave-absorption})
plug their hardware into. Most of this chapter is detector-agnostic;
specialisations to superconducting hardware appear as worked
numerical examples and are taken up systematically in the
subsequent chapters.

% --------------------------------------------------------------
\section{The dark-matter halo}
\label{sec:dm-halo}

\begin{phenomenon}
The Earth orbits within a roughly virialised, dynamically cold halo
of dark matter. Whatever else dark matter is, it has a definite
local density and a definite velocity distribution in the
laboratory frame, both of which are inputs to every direct-detection
rate calculation.
\end{phenomenon}

Several independent observations --- galactic rotation curves, the
peaks of the cosmic-microwave-background power spectrum,
gravitational lensing, and the abundance of light elements from
big-bang nucleosynthesis --- agree on a universe whose mass-energy
budget is roughly 5\% baryonic matter, 27\% non-luminous dark
matter, and 68\% dark energy~\cite{bertone2005}. The dark
component behaves gravitationally as cold, collisionless,
non-relativistic, non-baryonic matter; the search for the underlying
particle physics is what motivates the present chapter.

\begin{statement}[Standard halo model]
\label{stmt:halo}
The local dark-matter density is taken to be
\begin{equation}
\label{eq:rho-chi}
\rho_\chi \;\approx\; 0.3\;\text{GeV}/\text{cm}^3,
\end{equation}
and the velocity distribution in the galactic rest frame to be
isotropic Maxwell--Boltzmann, truncated at the galactic escape
velocity $v_\text{esc}$:
\begin{equation}
\label{eq:halo-distribution}
f_\text{gal}(\bm v) \;=\; \frac{1}{N_\text{esc}\,(\pi v_0^2)^{3/2}}\,
\exp\!\left(-\frac{|\bm v|^2}{v_0^2}\right)\,\Theta\!\left(v_\text{esc}-|\bm v|\right),
\end{equation}
with most-probable speed $v_0\approx 220\,\text{km}/\text{s}$ and
$v_\text{esc}\approx 544\,\text{km}/\text{s}$, normalised such that
$\int f_\text{gal}\,\dd^3v=1$.
\end{statement}

\begin{remark}[Lab-frame distribution]
The Earth moves through the halo at
$\bm v_E$ with $|\bm v_E|\approx 232\,\text{km}/\text{s}$ (the
solar-system motion plus the small annual modulation due to the
Earth's orbit). The lab-frame distribution is obtained by Galilean
boost,
$f_\text{lab}(\bm v)=f_\text{gal}(\bm v+\bm v_E)$. This is the
distribution that enters every rate formula below. The annual
modulation in $|\bm v_E|$ at the few-percent level is a powerful
discriminator of a putative dark-matter signal against
backgrounds, but plays no role in the leading-order analyses of
this chapter.
\end{remark}

The number density and flux follow immediately:
\begin{equation}
\label{eq:flux}
n_\chi \;=\; \frac{\rho_\chi}{m_\chi},
\qquad
\Phi_\chi \;=\; n_\chi\,\langle v\rangle,
\end{equation}
with $\langle v\rangle\sim v_0$ a velocity-averaged speed of order
the galactic dispersion. For $m_\chi=1\,\text{GeV}$, $\Phi_\chi\sim
10^7\,\text{cm}^{-2}\,\text{s}^{-1}$; for $m_\chi=1\,\text{MeV}$,
$\Phi_\chi\sim 10^{10}\,\text{cm}^{-2}\,\text{s}^{-1}$. Light DM is
plentiful but soft; heavy DM is rare but energetic. The crossover
that determines which of the next two channels is in business is
controlled by the detector threshold.

\paragraph{Transition.}
With local density and velocity in hand, we now treat the two
channels in turn: fermion-like DM that scatters
(\cref{sec:fermion-dm}) and bosonic DM that is absorbed
(\cref{sec:boson-dm}).

% --------------------------------------------------------------
\section{Fermion-like dark matter: detection by elastic recoil}
\label{sec:fermion-dm}

\begin{phenomenon}
A dark-matter particle that interacts via a contact-like
short-range force scatters elastically off a target particle in the
detector. The energy deposited is the recoil energy of the target
nucleus or electron; the cross-section per target encodes the
microphysics of the interaction.
\end{phenomenon}

% ----------- 6.3.1 Kinematics
\subsection{Kinematics: the recoil energy}
\label{sec:recoil-kinematics}

For a non-relativistic DM particle of mass $m_\chi$ and velocity
$\bm v$ scattering elastically off a target of mass $M_T$ at rest,
energy and momentum conservation in the lab frame fix the recoil
energy entirely in terms of the scattering angle.

\begin{statement}[Recoil-energy range]
\label{stmt:recoil-energy}
The recoil energy $E_R$ of the target after elastic scattering
runs from $0$ to a maximum
\begin{equation}
\label{eq:E_R-max}
E_R^\text{max} \;=\; \frac{2\,\mu_{\chi T}^2\,v^2}{M_T},
\qquad
\mu_{\chi T}\equiv\frac{m_\chi M_T}{m_\chi+M_T},
\end{equation}
where $\mu_{\chi T}$ is the reduced mass of the $\chi$--target
system. For a fixed recoil energy $E_R$, the minimum DM speed
needed to produce it is
\begin{equation}
\label{eq:v-min}
v_\text{min}(E_R) \;=\; \sqrt{\frac{M_T\,E_R}{2\,\mu_{\chi T}^2}}.
\end{equation}
\end{statement}

\begin{proof}
Take the lab frame in which the target is at rest. After
scattering, the DM and target have momenta $\bm p_\chi'$ and
$\bm p_T'$ satisfying
$\bm p_\chi'+\bm p_T'=m_\chi\bm v$ and
$|\bm p_\chi'|^2/(2m_\chi)+|\bm p_T'|^2/(2 M_T)=\tfrac12 m_\chi v^2$.
Eliminate $\bm p_\chi'$ between the two and use
$E_R=|\bm p_T'|^2/(2M_T)$ to obtain a quadratic in
$|\bm p_T'|$ whose physical root gives
$E_R(\theta)=2\mu_{\chi T}^2 v^2(1-\cos\theta^*)/M_T$ in the
centre-of-mass scattering angle $\theta^*$. The maximum
\cref{eq:E_R-max} corresponds to backscatter
($\cos\theta^*=-1$); the minimum velocity \cref{eq:v-min} is the
inverse relation $v(E_R)$ at $\theta^*=\pi$.
\end{proof}

\begin{example}[Light WIMP on a silicon nucleus]
\label{ex:wimp-on-Si}
For a $m_\chi=1\,\text{GeV}/c^2$ WIMP scattering off a silicon
nucleus ($M_T\approx 28\,m_p\approx 26\,\text{GeV}/c^2$), the
reduced mass is $\mu_{\chi T}\approx 0.96\,\text{GeV}/c^2$. With
$v\approx 220\,\text{km}/\text{s}=7.3\times 10^{-4}\,c$,
$E_R^\text{max}\approx 2(0.96)^2(7.3\times 10^{-4})^2/26\,\text{GeV}
\approx 38\,\text{eV}$. The maximum recoil falls below an eV for
$m_\chi\lesssim 100\,\text{MeV}$, a strong push to lower the
detector threshold.
\end{example}

% ----------- 6.3.2 Cross-section
\subsection{Differential cross-section}
\label{sec:recoil-xsec}

The cross-section per target encodes the unknown microphysics. We
quote the standard parametrisations used throughout the
direct-detection literature~\cite{lewin1996}, deferring derivations
to the references.

\begin{statement}[Spin-independent nuclear recoil]
\label{stmt:si-xsec}
For coherent (spin-independent) elastic scattering on a nucleus of
mass number $A$, the differential cross-section is
\begin{equation}
\label{eq:dsigma-dE-SI}
\frac{\dd\sigma_{\chi T}}{\dd E_R}
\;=\; \frac{M_T\,\sigma_n}{2\,\mu_{\chi n}^2\,v^2}\;A^2\,F^2(q),
\end{equation}
where $\sigma_n$ is the per-nucleon cross-section,
$\mu_{\chi n}=m_\chi m_n/(m_\chi+m_n)$ is the DM--nucleon reduced
mass, $A^2$ is the coherent enhancement, $q=\sqrt{2\,M_T\,E_R}$ is
the momentum transfer, and $F(q)$ is the nuclear form factor (we
use the standard Helm parametrisation~\cite{helm1956}, $F(0)=1$,
$F(q)\to 0$ at $qR\sim 1$ for nuclear radius $R$).
\end{statement}

\begin{remark}[Light vs heavy mediator]
\Cref{eq:dsigma-dE-SI} corresponds to a heavy-mediator (contact)
interaction; for a mediator mass $m_M\lesssim q$, the cross-section
acquires an extra propagator factor $\propto 1/(q^2+m_M^2)^2$, which
softens the recoil spectrum by suppressing large momentum transfers
and enhancing the small-$q$ region. We absorb this dependence into
the dimensionless DM form factor $F_\text{DM}(q)$, defined so that
$F_\text{DM}=1$ for a heavy mediator and
$F_\text{DM}(q)=(\alpha m_e/q)^2$ for an ultralight (massless or
sub-eV) mediator coupling to electrons through the photon.
\end{remark}

\paragraph{Transition.}
The nuclear-recoil cross-section above presumes that the target
particle is approximately free and recoils elastically. For a DM
particle scattering off a bound electron in a condensed-matter
target the recoil is intrinsically inelastic: energy goes into
ionising or exciting the electron system, and the kinematics are
no longer the constrained one-parameter family of the nuclear case.
The natural rewriting that keeps the formulas compact uses the
target's dielectric response.

\begin{statement}[DM--electron scattering: energy-loss function form]
\label{stmt:dme-xsec}
For DM scattering inelastically off the bound electrons of a
condensed-matter target, the momentum transfer $q$ and the energy
deposit $\omega$ (the electronic recoil energy) are independent
kinematic variables. The minimum DM speed required to transfer
$(q,\omega)$ is
\begin{equation}
\label{eq:vmin-electron}
v_\text{min}(q,\omega) \;=\; \frac{\omega}{q} \,+\, \frac{q}{2\,m_\chi}.
\end{equation}
The standard parametrisation of the inclusive scattering
cross-section per unit detector mass uses the target's longitudinal
\emph{energy-loss function} (ELF)
\begin{equation}
\label{eq:elf-def}
\mathcal{W}(q,\omega) \;\equiv\; -\,\text{Im}\!\left[\frac{1}{\hat\epsilon_L(q,\omega)}\right],
\end{equation}
where $\hat\epsilon_L(q,\omega)$ is the longitudinal complex
dielectric function. The differential rate per unit detector mass
in terms of $\mathcal{W}$ is
\begin{equation}
\label{eq:dRdE-electron}
\frac{1}{\mathcal{M}}\,\frac{\dd R}{\dd\omega}
\;=\; \frac{\rho_\chi}{m_\chi}\,
\frac{\bar\sigma_e}{\mu_{\chi e}^2}\,
\frac{1}{4\pi\,\alpha\,m_e^2\,\rho_T}\,
\int_{q_-(\omega)}^{q_+(\omega)}\!\! q^3\,\dd q\;
|F_\text{DM}(q)|^2\,\mathcal{W}(q,\omega)\,\eta\!\bigl(v_\text{min}(q,\omega)\bigr).
\end{equation}
Here $\bar\sigma_e$ is the reference DM--electron cross-section,
$\mu_{\chi e}=m_\chi m_e/(m_\chi+m_e)$ the DM--electron reduced
mass, $\rho_T$ the target mass density, $\alpha\approx 1/137$ the
fine-structure constant, $F_\text{DM}(q)$ the DM form factor of
the previous remark, $\eta$ the mean-inverse-speed function
\cref{eq:eta}, and $q_\pm(\omega)$ the kinematic limits at which
$v_\text{min}(q,\omega)=v_\text{esc}$.
\end{statement}

\paragraph{Why the ELF.}
The transition from the elastic, free-target picture of
\cref{stmt:si-xsec} to the bound, inelastic picture of
\cref{stmt:dme-xsec} amounts to replacing the nuclear form factor
$A^2 F^2(q)$ (a static, kinematic-only object) with the dynamical
response function $\mathcal{W}(q,\omega)$, which captures
band-structure, plasmon dispersion and screening of the target. The
fluctuation--dissipation theorem ensures that $\mathcal{W}$ is the
same object that controls the absorption of an external
electromagnetic perturbation at $(q,\omega)$, so the ELF is also
what one would measure by inelastic electron-energy-loss
spectroscopy on the same sample. This identification, and the
practical importance of having $\mathcal{W}$ tabulated as a
function of both arguments, is what motivates the numerical-tools
remark below.

\begin{remark}[Numerical evaluation of the energy-loss function]
The energy-loss function $\mathcal{W}(q,\omega)$ is, unlike the
nuclear form factor, intrinsically a many-body object: closed
forms exist only for toy models (free-electron gas, isolated
hydrogen-like atom). For real semiconductor and metal targets it
must be computed numerically from a band-structure calculation
(typically density-functional theory) or extrapolated from
optical-conductivity data. Public codes that tabulate
$\mathcal{W}(q,\omega)$ for the materials commonly used in
direct-detection experiments --- silicon, germanium, gallium
arsenide, aluminium and sapphire --- and that wrap the rate
integral \cref{eq:dRdE-electron} for use in analyses are
available~\cite{knapen2021}. Their numerical accuracy is limited by
the underlying band-structure calculation; in regions where
independent optical-loss data exist, agreement is typically at the
$\mathcal O(10\,\%)$ level.
\end{remark}

% ----------- 6.3.3 Predicted spectrum
\subsection{Predicted differential rate}
\label{sec:recoil-rate}

\begin{phenomenon}
The detector sees not a single recoil but a distribution of
recoils sampled over the velocity distribution of the DM halo.
The shape of the spectrum is the central observable of every
recoil search. Two case-specific formulas are required: one for
elastic nuclear recoils, where the recoil energy $E_R$ is the
single kinematic variable and the differential rate is
$\dd R/\dd E_R$; and one for inelastic electronic recoils, where
the deposited energy $\omega$ and the momentum transfer $q$ are
independent and the differential rate is $\dd R/\dd \omega$ at
fixed $q$, integrated over $q$.
\end{phenomenon}

\paragraph{Nuclear recoils.}
The nuclear case is the elastic free-target picture: $E_R$ and
$q$ are related kinematically by $q=\sqrt{2 M_T E_R}$, so the
velocity integral against the differential cross-section
\cref{eq:dsigma-dE-SI} is one-dimensional.

\begin{statement}[Velocity-averaged differential rate, nuclear recoils]
\label{stmt:dRdE}
For a target of total mass $\mathcal{M}$ comprised of nuclei of
individual mass $M_T$, the rate of elastic nuclear recoils per
unit detector mass and unit recoil energy is
\begin{equation}
\label{eq:dRdE_R}
\frac{1}{\mathcal M}\,\frac{\dd R}{\dd E_R}
\;=\; \frac{n_\chi}{m_T}\;\int_{v_\text{min}(E_R)}^{v_\text{esc}}\!\!
v\,f_\text{lab}(\bm v)\,\frac{\dd\sigma_{\chi T}}{\dd E_R}\,\dd^3v,
\end{equation}
with $v_\text{min}(E_R)$ given by \cref{eq:v-min} and
$\dd\sigma_{\chi T}/\dd E_R$ by \cref{eq:dsigma-dE-SI}.
\end{statement}

\begin{proof}
Write the rate at which DM particles in a velocity bin
$[\bm v,\bm v+\dd^3 v]$ scatter into a recoil-energy bin
$[E_R,E_R+\dd E_R]$ as the product (number density of DM in that
velocity bin) $\times$ (number density of targets) $\times$
(relative speed) $\times$ (differential cross-section). Each is
read off from \cref{eq:flux}, \cref{eq:halo-distribution} and
\cref{eq:dsigma-dE-SI}; the $v_\text{min}$ lower limit on the
velocity integral expresses the kinematic requirement
\cref{eq:v-min}.
\end{proof}

\paragraph{Mean-inverse-speed function (nuclear case).}
For the standard halo model and the elastic-recoil kinematics
above, the velocity integral in \cref{eq:dRdE_R} is conventionally
written through the mean-inverse-speed function
\begin{equation}
\label{eq:eta}
\eta(v_\text{min}) \;\equiv\; \int_{v\geq v_\text{min}}\frac{f_\text{lab}(\bm v)}{v}\,\dd^3v,
\end{equation}
which has a closed form in terms of error functions (see
e.g.~\cite{lewin1996} for the explicit expression). Plugging
\cref{eq:dsigma-dE-SI} into \cref{eq:dRdE_R}, the differential
rate takes the often-quoted form
\begin{equation}
\label{eq:dRdE-final}
\frac{1}{\mathcal M}\,\frac{\dd R}{\dd E_R}
\;=\; \frac{\rho_\chi\,\sigma_n}{2\,m_\chi\,\mu_{\chi n}^2}\;A^2\,F^2(q)\;\eta(v_\text{min}(E_R)).
\end{equation}

\begin{example}[Nuclear-recoil spectrum on silicon]
\label{ex:Si-spectrum}
For $m_\chi=10\,\text{GeV}$, $\sigma_n=10^{-40}\,\text{cm}^2$, on
silicon ($A=28$): \cref{eq:dRdE-final} gives a spectrum that peaks
at $E_R\approx 0$ and falls off exponentially with characteristic
scale $E_R^\text{max}\approx 1\,\text{keV}$. A detector with
threshold $E_\text{th}\sim 100\,\text{eV}$ retains the bulk of the
spectrum; $E_\text{th}\sim 1\,\text{keV}$ retains essentially none
of it. This trade-off, between threshold and dark-matter mass
reach, is what motivates the cryogenic detection programme.
\end{example}

\paragraph{Electron recoils.}
The electron case is qualitatively different because the
electron is bound to the nucleus and the recoil is inelastic on
the atom or crystal: the momentum transfer $q$ and the deposited
energy $\omega$ are independent kinematic variables, and the
mean-inverse-speed function $\eta(v_\text{min}(E_R))$ used above
is replaced by $\eta(v_\text{min}(q,\omega))$ inside an integral
over $q$.

\begin{statement}[Velocity-averaged differential rate, electron recoils]
\label{stmt:dRdomega}
For a target of total mass $\mathcal M$ and bulk mass density
$\rho_T$, the rate of inelastic DM-electron events per unit
detector mass and unit deposited energy is
\begin{equation}
\label{eq:dRdomega-restated}
\frac{1}{\mathcal M}\,\frac{\dd R}{\dd \omega}
\;=\; \frac{\rho_\chi}{m_\chi}\,
\frac{\bar\sigma_e}{\mu_{\chi e}^2}\,
\frac{1}{4\pi\,\alpha\,m_e^2\,\rho_T}\,
\int_{q_-(\omega)}^{q_+(\omega)}\!\! q^3\,\dd q\;
|F_\text{DM}(q)|^2\,\mathcal{W}(q,\omega)\,
\eta\!\bigl(v_\text{min}(q,\omega)\bigr),
\end{equation}
with $v_\text{min}(q,\omega)$ given by \cref{eq:vmin-electron},
$\mathcal{W}(q,\omega)$ the energy-loss function of
\cref{eq:elf-def}, and $q_\pm(\omega)$ the kinematic limits at
which $v_\text{min}(q,\omega)=v_\text{esc}$.
\Cref{eq:dRdomega-restated} is identical to \cref{eq:dRdE-electron};
we restate it here so that \cref{stmt:dRdE} and \cref{stmt:dRdomega}
sit side by side.
\end{statement}

\begin{proof}
Repeat the bin-counting argument of \cref{stmt:dRdE} with the
DM--electron differential cross-section in place of
$\dd\sigma_{\chi T}/\dd E_R$. The inelasticity of the recoil
turns the one-parameter family $E_R\leftrightarrow q$ of the
nuclear case into the two-parameter family $(q,\omega)$, so that
the differential rate at fixed $\omega$ is itself a $q$-integral.
The role played by $\dd\sigma_{\chi T}/\dd E_R\,A^2 F^2(q)$ in the
nuclear case is taken by the energy-loss function $\mathcal{W}(q,\omega)$
weighted by the DM form factor $|F_\text{DM}(q)|^2$; the velocity
integral for each $(q,\omega)$ pair gives $\eta(v_\text{min}(q,\omega))$.
\end{proof}

\begin{example}[Electron-recoil spectrum on silicon]
\label{ex:Si-electron-spectrum}
For $m_\chi=10\,\text{MeV}$, $\bar\sigma_e=10^{-37}\,\text{cm}^2$,
a heavy mediator ($F_\text{DM}=1$), and a silicon target
($\rho_T=2.33\,\text{g/cm}^3$), \cref{eq:dRdomega-restated}
evaluated with the publicly tabulated~\cite{knapen2021} ELF for
crystalline silicon yields a spectrum dominated by the
bulk-plasmon resonance near $\omega\approx 17\,\text{eV}$, with a
low-energy tail above the indirect band gap of
$E_g\approx 1.1\,\text{eV}$. The characteristic kinematic threshold
is set by demanding that
$v_\text{min}(q,\omega)\leq v_\text{esc}$ in
\cref{eq:vmin-electron}; with $\omega/q\sim 10^{-3}\,c$ around the
gap, this places the lowest detectable DM mass at
$m_\chi^\text{min}\sim E_g/(2 v_\text{esc}^2 m_e)\,m_e\approx
0.5\,\text{MeV}$ for $\omega=E_g$. A detector with
$E_\text{th}=E_g$ thus reaches roughly two orders of magnitude
lower in DM mass than a comparable nuclear-recoil analysis on the
same crystal.
\end{example}

\paragraph{Transition.}
This is the predicted in-detector spectrum for the recoil channel.
Before turning to the detector-response and statistics steps that
convert it into a measurement and a limit, we work the parallel
walk-through for the absorption channel.

% --------------------------------------------------------------
\section{Bosonic dark matter: detection by absorption}
\label{sec:boson-dm}

\begin{phenomenon}
A sufficiently light bosonic DM candidate (mass below an eV) is
better thought of as a coherent classical field filling the
detector volume than as a stream of particles. If the field couples
to ordinary matter, the detector can absorb a quantum from it,
depositing the full rest energy $m_\chi c^2$ at once. The signature
is a narrow, near-monochromatic line at $f_\chi=m_\chi c^2/h$,
broadened only by the small virial dispersion of the halo.
\end{phenomenon}

% ----------- 6.4.1 Kinematics of absorption
\subsection{Kinematics of absorption: a virial-broadened line}
\label{sec:absorption-kinematics}

\begin{statement}[Line spectrum at the Compton frequency]
\label{stmt:line}
Absorption of a non-relativistic bosonic DM particle deposits
energy
\begin{equation}
\label{eq:E-deposited}
E_\text{dep} \;=\; m_\chi c^2 \,+\, \tfrac12 m_\chi v^2
\;\approx\; m_\chi c^2\,\bigl(1+\tfrac{1}{2}(v/c)^2\bigr).
\end{equation}
The kinetic correction is small: with $v/c\sim 10^{-3}$ the line
sits at $f_\chi=m_\chi c^2/h$ with fractional width
$\Delta f/f_\chi \sim (v/c)^2/2 \sim 10^{-6}$.
\end{statement}

\begin{example}[Frequency--mass dictionary]
\label{ex:freq-mass}
Useful conversions:
$1\,\mu\text{eV}\leftrightarrow 242\,\text{MHz}$,
$4\,\mu\text{eV}\leftrightarrow 1\,\text{GHz}$,
$1\,\text{meV}\leftrightarrow 242\,\text{GHz}$,
$1\,\text{eV}\leftrightarrow 242\,\text{THz}$. The QCD-axion
window of $1\text{--}40\,\mu\text{eV}$ corresponds to
$240\,\text{MHz}\text{--}10\,\text{GHz}$, exactly the band in
which superconducting transmons operate.
\end{example}

% ----------- 6.4.2 Coupling Lagrangians
\subsection{Coupling Lagrangians: dark photons and axions}
\label{sec:absorption-couplings}

The two photon-mediated bosonic candidates that dominate the
literature, and the only ones we treat here, are the dark photon
$A'_\mu$ and the QCD axion $a$. Each couples to ordinary matter
through a single dimensionful parameter --- the kinetic-mixing
parameter $\varepsilon$ for dark photons, or the photon coupling
$g_{a\gamma\gamma}$ for axions --- and the detection rate is
quadratic in that parameter.

\begin{statement}[Dark photon kinetic mixing]
\label{stmt:dark-photon}
A dark photon $A'_\mu$ kinetically mixes with the Standard-Model
photon~\cite{holdom1986},
\begin{equation}
\label{eq:dark-photon-L}
\mathcal{L}\;\supset\;-\tfrac{1}{4}F'_{\mu\nu}F'^{\mu\nu}-\tfrac12 m_{A'}^2 A'_\mu A'^\mu
+\tfrac{\varepsilon}{2}F'_{\mu\nu}F^{\mu\nu}.
\end{equation}
A field redefinition transfers the mixing into the matter sector,
where the dark photon couples to the electromagnetic current with
effective charge $\varepsilon e$. A mass eigenstate $A'$ identified
with cold dark matter has, in any sample of volume $V$, an average
field amplitude $\langle |\bm A'|^2\rangle\propto \rho_\chi/m_{A'}^2$
oscillating at the Compton frequency $\omega_{A'}=m_{A'}c^2/\hbar$.
\end{statement}

\begin{statement}[Axion-photon coupling]
\label{stmt:axion}
The axion couples to two photons through
\begin{equation}
\label{eq:axion-L}
\mathcal{L}_{a\gamma\gamma} \;=\; -\tfrac14\,g_{a\gamma\gamma}\,a\,F_{\mu\nu}\tilde F^{\mu\nu}
\;=\; g_{a\gamma\gamma}\,a\,\bm{E}\!\cdot\!\bm{B}.
\end{equation}
In an external static magnetic field $\bm B_0$, the axion field
$a(t)\sim\sqrt{2\rho_\chi}/m_a\,\cos(m_a c^2 t/\hbar)$ acts as a
time-dependent source for an oscillating electric field
$\bm E_a\sim g_{a\gamma\gamma}\,a\,\bm B_0$. Inside a microwave
cavity tuned to $\omega_a=m_a c^2/\hbar$, $\bm E_a$ drives the
resonant TM mode coherently; this is the Sikivie haloscope
concept~\cite{sikivie1983}.
\end{statement}

% ----------- 6.4.3 Absorption rate
\subsection{Predicted absorption rate}
\label{sec:absorption-rate}

\begin{phenomenon}
A coherent oscillating DM field acts on the detector as an external
classical drive at frequency $m_\chi c^2/\hbar$. By Fermi's golden
rule, the absorption rate per absorbing target is set by the matrix
element of the coupling and by the detector's response at that
frequency.
\end{phenomenon}

\begin{statement}[Dark-photon absorption rate]
\label{stmt:dark-photon-rate}
For a dark photon with mass $m_{A'}$ and kinetic mixing $\varepsilon$,
the absorption rate per unit absorber volume is
\begin{equation}
\label{eq:dark-photon-rate}
\frac{\dd N_\text{abs}}{\dd t \,\dd V} \;=\; \varepsilon^2\,\frac{\rho_\chi}{m_{A'}}\,
\sigma_\text{abs}^\text{ord}(\omega_{A'}),
\end{equation}
where $\sigma_\text{abs}^\text{ord}(\omega)$ is the ordinary-photon
absorption cross-section per unit volume of the same absorber at
frequency $\omega$, and the right-hand side is evaluated at
$\omega=\omega_{A'}=m_{A'}c^2/\hbar$.
\end{statement}

\paragraph{Derivation.}
The kinetic mixing in \cref{eq:dark-photon-L} makes the dark
photon, when present at density $\rho_\chi$, equivalent to an
ordinary photon flux of suppression $\varepsilon^2$ at frequency
$\omega_{A'}$~\cite{caputo2021review}. The probability per unit
time per unit volume of absorbing one such ordinary photon is
$\Phi_\text{ord}\,\sigma_\text{abs}^\text{ord}$, with $\Phi_\text{ord}$
the equivalent photon flux. Substituting
$\Phi_\text{ord}=\varepsilon^2\,\rho_\chi/m_{A'}$ (the local DM
energy density rewritten as a photon flux) gives
\cref{eq:dark-photon-rate}. The same logic applied via Fermi's
golden rule (\cref{thm:fermi}), with $\op V$ the kinetically-mixed
electromagnetic interaction, yields the same answer.

\begin{statement}[Axion haloscope signal power]
\label{stmt:axion-power}
For an axion of mass $m_a$ and coupling $g_{a\gamma\gamma}$, an
external magnetic field $B_0$, and a microwave cavity of volume
$V$ and loaded quality factor $Q_L$ tuned to $\omega_a=m_a c^2/\hbar$,
the on-resonance signal power deposited in the cavity is
\begin{equation}
\label{eq:axion-power}
P_\text{sig} \;=\; \kappa\,g_{a\gamma\gamma}^2\,\rho_\chi\,
\frac{B_0^2 V}{m_a^2}\,\omega_a\,Q_L\,C_{nlm},
\end{equation}
with $C_{nlm}\in[0,1]$ the form-factor of the cavity mode and
$\kappa\in[0,1]$ the cavity-coupling factor (typically $\kappa=1/2$
for critical coupling).
\end{statement}

\begin{example}[Numerical absorption rate, dark photon at 1 meV]
\label{ex:dark-photon-rate}
Take $m_{A'}=1\,\text{meV}/c^2$ ($\omega_{A'}/2\pi=242\,\text{GHz}$),
$\rho_\chi=0.3\,\text{GeV}/\text{cm}^3$, and an aluminium absorber
of volume $1\,\text{cm}^3$ with $\sigma_\text{abs}^\text{ord}\sim
10^4\,\text{cm}^{-1}$ at $1\,\text{meV}$ (well above the
$2\Delta\approx 0.34\,\text{meV}$ Cooper-pair threshold). For
$\varepsilon=10^{-13}$, \cref{eq:dark-photon-rate} gives an
absorption rate of order $10^{-3}\,\text{Hz}/\text{cm}^3$, or
$\sim 10^2$ counts per kg per day --- comparable with the
quasiparticle background of present-generation devices.
\end{example}

\begin{example}[Numerical signal power, 4 µeV axion haloscope]
\label{ex:axion-power}
For $m_a=4\,\mu\text{eV}/c^2$ ($\omega_a/2\pi=1\,\text{GHz}$),
$g_{a\gamma\gamma}=10^{-15}\,\text{GeV}^{-1}$,
$B_0=8\,\text{T}$, $V=1\,\text{L}$, $Q_L=10^5$, $C_{010}\approx 0.7$,
$\kappa=1/2$, \cref{eq:axion-power} gives
$P_\text{sig}\sim 10^{-23}\,\text{W}\sim 10^{-2}\,\text{photons}/\text{s}$
in a cavity bandwidth, marginally detectable with a
single-microwave-photon counter (the resonant scan covers a
relative-bandwidth $\sim 1/Q_L$ per integration time).
\end{example}

\paragraph{Transition.}
Both channels now have a predicted in-detector spectrum
($\dd R/\dd E_R$ for recoils, a near-monochromatic line at
$\omega=m_\chi c^2/\hbar$ for absorption). The next two sections
turn that spectrum into a number of detector counts and then into
a statistical limit on the underlying coupling.

% --------------------------------------------------------------
\section{From spectrum to detected counts}
\label{sec:detector-response}

\begin{phenomenon}
A real detector samples the predicted spectrum with finite energy
threshold, finite efficiency, and finite resolution. The measurable
quantity is not the underlying signal spectrum but its convolution
with the detector response, integrated over the experimental
exposure.
\end{phenomenon}

\begin{statement}[Expected event count]
\label{stmt:N-exp}
For a detector with exposure $\mathcal E$ (target mass times live
time), efficiency $\eta(E)$, and Gaussian energy-resolution kernel
\begin{equation}
G_\sigma(E\,|\,E') \;=\; \frac{1}{\sqrt{2\pi}\,\sigma(E')}\,
\exp\!\left(-\frac{(E-E')^2}{2\sigma^2(E')}\right),
\end{equation}
the expected number of signal events above threshold $E_\text{th}$
is
\begin{equation}
\label{eq:N-exp}
N_\text{exp} \;=\; \mathcal E \int_{E_\text{th}}^\infty\!\!\dd E\;
\eta(E)\,\bigl[(\dd R/\dd E')\!\ast\!G_\sigma\bigr](E),
\end{equation}
where the convolution acts on the predicted differential rate
($\dd R/\dd E_R$ for recoils, or a Gaussian of width
$\Delta f/f\sim 10^{-6}$ centred at $f_\chi$ for absorption).
\end{statement}

\begin{remark}[Resonant absorption searches]
For absorption searches with a resonant cavity (the haloscope),
``threshold and resolution'' are replaced by the cavity
bandwidth $\Delta\omega\sim\omega/Q_L$ and the scan strategy: each
cavity tune covers a frequency range $\sim\omega/Q_L$ for a time
$\tau$, and the experimentally relevant quantity is the integrated
power $P_\text{sig}\,\tau$ collected per scan step.
\end{remark}

\begin{example}[Phonon detector, recoil channel]
\label{ex:phonon-detector}
A $1\,\text{kg}$ silicon target with athermal-phonon readout has
$E_\text{th}\sim 10\,\text{meV}$, $\sigma(E_R)\sim 1\,\text{meV}$,
and $\eta(E_R)\approx 1$ above threshold. With a $30$-day live
time, $\mathcal E=30\,\text{kg-day}$. For
$m_\chi=100\,\text{MeV}$, $\sigma_n=10^{-37}\,\text{cm}^2$ and the
spectrum of \cref{eq:dRdE-final}, \cref{eq:N-exp} gives
$N_\text{exp}\sim 10$ events --- at the threshold of detectability,
just above background.
\end{example}

\begin{example}[Transmon counter, absorption channel]
\label{ex:transmon-counter}
A transmon at $\omega_q/2\pi=5\,\text{GHz}$ tuned to $\omega_{A'}$
acts as a single-microwave-photon detector with photon-counting
efficiency $\eta\sim 0.5$ and dark-count rate
$\Gamma_\text{dc}\sim 1\,\text{mHz}$~\cite{dixit2021}. For
$\varepsilon=10^{-15}$ and an integration time of $1\,\text{day}$,
\cref{eq:dark-photon-rate} integrated over the cavity bandwidth
gives $\langle N_\text{exp}\rangle\sim 1$, comparable with the
expected number of dark counts ($\Gamma_\text{dc}\,\tau\sim 100$
in a day, but the resonant cavity bandwidth narrows this
substantially).
\end{example}

% --------------------------------------------------------------
\section{From counts to limits}
\label{sec:limits}

\begin{phenomenon}
A finite exposure with $N$ observed events constrains the
underlying coupling. In the absence of an observed signal one
quotes an upper limit; in the presence of one, a confidence
interval. Both reduce, in the simplest cases, to thresholding a
likelihood function constructed from a Poisson model of the count.
\end{phenomenon}

\subsection{Poisson upper limit}
\label{sec:poisson-limit}

\begin{statement}[Poisson 90 \% CL upper limit]
\label{stmt:poisson-UL}
Suppose the expected count is the sum of a signal contribution
$\mu(\theta)$ depending on the coupling parameter $\theta$ (with
$\mu\propto\theta^2$ for the bosonic and $\mu\propto\sigma_n$ for
the fermionic case) and a known background rate $b$. The number of
observed events $N$ is Poisson-distributed with mean
$\mu(\theta)+b$. The 90\% confidence-level upper limit
$\theta_{90}$ on the coupling is the smallest $\theta$ such that
\begin{equation}
\label{eq:poisson-UL}
P(n\leq N \,|\,\mu(\theta_{90})+b)
\;=\;\sum_{n=0}^{N}\frac{[\mu(\theta_{90})+b]^n\,e^{-[\mu(\theta_{90})+b]}}{n!}
\;\leq\; 0.10.
\end{equation}
For $b=0$ and $N=0$, $\mu_{90}=2.30$; for $N=1$, $\mu_{90}=3.89$;
for $N=2$, $\mu_{90}=5.32$.
\end{statement}

\begin{remark}[Feldman--Cousins for low counts]
The naive Poisson upper limit \cref{eq:poisson-UL} can over-cover
or under-cover the true coverage in the small-$N$ regime,
particularly when $N$ fluctuates below the expected background. The
\emph{Feldman--Cousins unified construction}~\cite{feldman1998}
provides a properly-covering interval that smoothly transitions
between two-sided and upper-limit regimes. We use the simpler
\cref{eq:poisson-UL} when $b\ll N_\text{exp}$ and Feldman--Cousins
otherwise.
\end{remark}

\subsection{Profile likelihood}
\label{sec:profile-likelihood}

\begin{phenomenon}
Real experiments have nuisance parameters --- detection efficiency,
background normalisation, energy calibration --- whose uncertainty
must be propagated into the limit. The profile-likelihood approach
maximises the likelihood over the nuisance parameters at each value
of the parameter of interest, producing a one-dimensional test
statistic.
\end{phenomenon}

\begin{statement}[Profile likelihood ratio]
\label{stmt:profile-likelihood}
Let $\bm{\theta}=(\mu,\bm{\nu})$ where $\mu$ is the parameter of
interest (signal strength) and $\bm{\nu}$ the nuisance parameters.
For an observed dataset $\bm x$, the likelihood is
$L(\mu,\bm\nu\,|\,\bm x)$ and the profile likelihood ratio is
\begin{equation}
\label{eq:profile-ratio}
\lambda(\mu) \;=\; \frac{L(\mu,\hat{\hat{\bm\nu}}(\mu)\,|\,\bm x)}
{L(\hat\mu,\hat{\bm\nu}\,|\,\bm x)},
\end{equation}
where $\hat\mu,\hat{\bm\nu}$ are the global maximum-likelihood
estimators and $\hat{\hat{\bm\nu}}(\mu)$ is the conditional maximum
of $\bm\nu$ at fixed $\mu$. Under the regularity conditions of
Wilks' theorem, the test statistic
\begin{equation}
\label{eq:wilks}
t_\mu \;\equiv\; -2\ln\lambda(\mu)
\end{equation}
follows a $\chi^2_1$ distribution asymptotically, and the
$1-\alpha$ confidence interval is the set of $\mu$ for which
$t_\mu\leq F^{-1}_{\chi^2_1}(1-\alpha)$ (e.g.\ $t_\mu\leq 2.71$ for
$1-\alpha=0.90$).
\end{statement}

\paragraph{Derivation.}
For Poisson data with mean $\mu+b$, the log-likelihood is
$\ln L=-\mu-b+N\ln(\mu+b)$, with the global maximum at
$\hat\mu=N-b$. Substituting in \cref{eq:wilks},
$t_\mu=2[(\mu-\hat\mu)+N\ln(\hat\mu/\mu)]$ at $b=0$, and the
$\chi^2_1$ asymptotic distribution follows from the standard
saddle-point expansion around the maximum (Cowan, Cranmer, Gross
and Vitells \cite{cowan2011} give explicit closed-form expressions
including the case of one or two-sided intervals).

\begin{example}[Comparison of the two limits, no nuisance parameters]
\label{ex:limits-compare}
For a single-bin Poisson experiment with $N=0$ observed events and
no background, the simple Poisson upper limit \cref{eq:poisson-UL}
is $\mu_{90}=2.30$. The corresponding profile-likelihood limit
\cref{eq:wilks} gives $t_\mu=2\mu$ (for $b=0$, $N=0$), so
$\mu_{90}=1.35$ for one-sided $90\%$ CL. The two differ because
the profile-likelihood interval uses a likelihood ratio rather
than a tail-probability cut; both are widely quoted in
direct-detection papers, and \cref{eq:poisson-UL} is the more
conservative.
\end{example}

\subsection{From the limit to an exclusion curve}
\label{sec:exclusion-curves}

The mass of the DM candidate enters the limit indirectly, through
the differential rate \cref{eq:dRdE-final} or absorption rate
\cref{eq:dark-photon-rate}. Solving \cref{eq:poisson-UL} or
\cref{eq:wilks} for the coupling at each mass produces a
one-dimensional curve --- the exclusion curve $\sigma_n(m_\chi)$,
$\varepsilon(m_{A'})$, or $g_{a\gamma\gamma}(m_a)$ --- that bounds
the model parameter space. The shape of the curve is dominated by
two effects: the kinematic threshold $E_R^\text{max}\geq E_\text{th}$
(in the recoil case) or the resonant-bandwidth coverage of the
detector (in the haloscope case), and the energy-dependent
detection efficiency $\eta(E)$.

\begin{example}[Sketch of an exclusion curve]
\label{ex:exclusion-sketch}
For the silicon phonon detector of \cref{ex:phonon-detector}, no
observed events at $E_\text{th}=10\,\text{meV}$ and exposure
$\mathcal E=30\,\text{kg-day}$ would set
$\mu_{90}=2.30$ on the expected count, translating to
$\sigma_n\lesssim 10^{-39}\,\text{cm}^2$ at $m_\chi=100\,\text{MeV}$.
The curve cuts off sharply below $m_\chi\sim 30\,\text{MeV}$ where
$E_R^\text{max}<E_\text{th}$.
\end{example}

% --------------------------------------------------------------
This chapter walked through two distinct dark-matter detection
channels --- elastic recoil for fermion-like DM and absorption for
bosonic DM --- and showed that the same four-stage pipeline (model
parameters, predicted spectrum, detected counts, statistical limit)
applies to both. The recoil channel produces a continuum spectrum
$\dd R/\dd E_R$ whose shape is set by the kinematic threshold
$E_R^\text{max}=2\mu_{\chi T}^2 v^2/M_T$ and the velocity-averaged
cross-section through the mean-inverse-speed function $\eta(v_\text{min})$.
The absorption channel produces a near-monochromatic line at
$\omega=m_\chi c^2/\hbar$ broadened only by the halo virial
dispersion to $\Delta\omega/\omega\sim 10^{-6}$. Either signal,
after folding with the detector response, becomes a Poisson count
whose expectation depends on the coupling parameter; thresholding a
likelihood produces an exclusion curve in the coupling-mass plane.

The next four chapters take this generic apparatus and instantiate
it on superconducting hardware: phonon-quasiparticle conversion
(\cref{ch:phonon-qp}) and piezoelectric phonon-to-qubit coupling
(\cref{ch:phonon-qubit-piezo}) for the recoil channel, and
infrared absorption in the Josephson junction
(\cref{ch:ir-junction}) and direct microwave absorption in the
transmon (\cref{ch:microwave-absorption}) for the absorption
channel. The methodology of this chapter applies to all four
uniformly.

\clearpage
\begin{chaptersummary}

\paragraph{Two detection channels.}
Direct-detection searches split cleanly into two cases. A
fermion-like DM particle scatters elastically off a target and
deposits a recoil energy; a bosonic DM field is absorbed wholesale
and deposits its rest energy $m_\chi c^2$ at once.

\paragraph{Halo inputs (\cref{sec:dm-halo}).}
Local density $\rho_\chi\approx 0.3\,\text{GeV}/\text{cm}^3$;
truncated isotropic Maxwellian velocity distribution
\cref{eq:halo-distribution} with $v_0\approx 220\,\text{km}/\text{s}$
and $v_\text{esc}\approx 544\,\text{km}/\text{s}$ (galactic frame);
flux $\Phi_\chi=n_\chi\bar v$.

\paragraph{Recoil pipeline (\cref{sec:fermion-dm}).}
\begin{itemize}
  \item Kinematics: $E_R^\text{max}=2\mu_{\chi T}^2 v^2/M_T$.
  \item Cross-section: spin-independent
        $\dd\sigma/\dd E_R=M_T\sigma_n A^2 F^2(q)/(2\mu_{\chi n}^2 v^2)$.
  \item Differential rate
        $\dd R/\dd E_R=\rho_\chi\sigma_n A^2 F^2(q)\,\eta(v_\text{min})/(2 m_\chi\mu_{\chi n}^2)$
        per unit detector mass.
\end{itemize}

\paragraph{Absorption pipeline (\cref{sec:boson-dm}).}
\begin{itemize}
  \item Line at $E=m_\chi c^2$ with fractional width
        $\Delta E/E\sim (v/c)^2/2\sim 10^{-6}$.
  \item Dark-photon rate per unit volume:
        $\dd N/(\dd t\,\dd V)=\varepsilon^2\rho_\chi\sigma_\text{abs}^\text{ord}/m_{A'}$.
  \item Axion-haloscope on-resonance signal power
        $P_\text{sig}=\kappa g_{a\gamma\gamma}^2\rho_\chi B_0^2 V \omega_a Q_L C_{nlm}/m_a^2$.
\end{itemize}

\paragraph{Detector response (\cref{sec:detector-response}).}
$N_\text{exp}=\mathcal E\!\int_{E_\text{th}}^\infty\!\dd E\,\eta(E)\,
\bigl[(\dd R/\dd E')\!\ast\!G_\sigma\bigr](E)$.

\paragraph{Statistical limits (\cref{sec:limits}).}
Poisson 90\%~CL upper limit
$\sum_{n=0}^N P(n|\mu_{90})\leq 0.10$ (e.g.\ $\mu_{90}=2.30$ for
$N=0$, $b=0$); profile-likelihood ratio
$\lambda(\mu)=L(\mu,\hat{\hat\nu}(\mu))/L(\hat\mu,\hat\nu)$ with
test statistic $t_\mu=-2\ln\lambda(\mu)\sim\chi^2_1$ asymptotically.

\paragraph{Hardware specialisation.}
The four mechanism-specific chapters that follow
(\cref{ch:phonon-qp,ch:ir-junction,ch:phonon-qubit-piezo,ch:microwave-absorption})
specialise this generic apparatus to superconducting hardware.

\end{chaptersummary}

% --------------------------------------------------------------
\section*{Exercises}
\addcontentsline{toc}{section}{Exercises}

The exercises grow in difficulty from inspection-level checks of
the formulas to applied calculations of the kind one runs when
designing a search.

\begin{exercise}[DM number density and flux]
For $\rho_\chi=0.3\,\text{GeV}/\text{cm}^3$ and a galactic
mean speed $\bar v\approx 240\,\text{km}/\text{s}$, compute the
number density $n_\chi$ and flux $\Phi_\chi$ for
$m_\chi\in\{1\,\text{MeV},\,1\,\text{GeV},\,100\,\text{GeV}\}$.
\begin{solution}
$n_\chi=\rho_\chi/m_\chi$ and $\Phi_\chi=n_\chi\bar v$. With
$\bar v=2.4\times 10^7\,\text{cm}/\text{s}$:
$n_\chi=\{300,\,0.3,\,3\times 10^{-3}\}\,\text{cm}^{-3}$ and
$\Phi_\chi=\{7\times 10^9,\,7\times 10^6,\,7\times 10^4\}\,\text{cm}^{-2}\text{s}^{-1}$.
Light dark matter is plentiful; heavy dark matter is rare.
\end{solution}
\end{exercise}

\begin{exercise}[Maximum recoil energy on silicon]
Compute $E_R^\text{max}$ for a silicon nucleus ($A=28$,
$M_T\approx 26\,\text{GeV}/c^2$) struck by a DM particle of mass
$m_\chi=1\,\text{GeV}/c^2$ and speed $v=220\,\text{km}/\text{s}$.
What if instead $m_\chi=10\,\text{MeV}/c^2$?
\begin{solution}
Reduced mass $\mu_{\chi T}=m_\chi M_T/(m_\chi+M_T)$.
For $m_\chi=1\,\text{GeV}$, $\mu_{\chi T}\approx 0.96\,\text{GeV}$,
$v/c=7.3\times 10^{-4}$, so
$E_R^\text{max}=2(0.96)^2(7.3\times 10^{-4})^2/26\,\text{GeV}\approx 38\,\text{eV}$.
For $m_\chi=10\,\text{MeV}$, $\mu_{\chi T}\approx 10\,\text{MeV}$,
$E_R^\text{max}\approx 2(10\times 10^{-3})^2(7.3\times 10^{-4})^2/26\,\text{GeV}
\approx 4\,\text{meV}$ --- below the meV threshold of any current
phonon detector.
\end{solution}
\end{exercise}

\begin{exercise}[Threshold-driven mass reach]
A phonon detector has threshold $E_\text{th}=10\,\text{meV}$ on
silicon. What is the minimum DM mass for which a recoil at
$v=v_\text{esc}\approx 544\,\text{km}/\text{s}$ can produce a
detectable event?
\begin{solution}
Demand $E_R^\text{max}(v_\text{esc})\geq E_\text{th}$, i.e.\
$2\mu_{\chi T}^2 v_\text{esc}^2/M_T\geq E_\text{th}$. Solving for
$\mu_{\chi T}^2$ with $v_\text{esc}/c=1.8\times 10^{-3}$ and
$M_T=26\,\text{GeV}$ gives
$\mu_{\chi T}^2\geq E_\text{th}M_T/(2 v_\text{esc}^2)
\approx 4\times 10^{-2}\,\text{GeV}^2$, i.e.\
$\mu_{\chi T}\gtrsim 200\,\text{MeV}$. Since
$\mu_{\chi T}\to m_\chi$ for $m_\chi\ll M_T$, this gives
$m_\chi\gtrsim 200\,\text{MeV}$. A 10\,meV phonon detector on
silicon thus probes nuclear recoils for $m_\chi$ above a few
hundred MeV.
\end{solution}
\end{exercise}

\begin{exercise}[Compton frequency of an axion]
The QCD-axion mass window is $1\,\mu\text{eV}\lesssim m_a c^2\lesssim
40\,\mu\text{eV}$. Compute the corresponding Compton frequency
range $f_a=m_a c^2/h$ and the fractional line width
$\Delta f/f_a$.
\begin{solution}
$f_a=m_a c^2/h$, with $1\,\mu\text{eV}=0.242\,\text{GHz}$, so
$f_a\in[0.24,\,9.7]\,\text{GHz}$. Fractional line width
$\Delta f/f_a\sim (v/c)^2/2\sim 10^{-6}$, i.e.\ a few hundred Hz to
a few kHz absolute width across the window.
\end{solution}
\end{exercise}

\begin{exercise}[Coherence volume of a dark photon field]
The bosonic-DM picture is valid when the de Broglie wavelength
$\lambda_\text{dB}\sim h/(m_\chi v)$ is much larger than the
detector. For a $1\,\mu\text{eV}$ dark photon at $v=220\,\text{km}/\text{s}$,
compute $\lambda_\text{dB}$ and compare it with a 1 cm cavity.
\begin{solution}
$p=m_\chi v=10^{-6}\,\text{eV}/c\cdot 7\times 10^{-4}\,c=7\times 10^{-10}\,\text{eV}/c$.
$\lambda_\text{dB}=h/p=(6.6\times 10^{-25}\,\text{eV s}\cdot
3\times 10^8\,\text{m/s})/(7\times 10^{-10}\,\text{eV})\approx 280\,\text{m}$.
Far larger than a 1 cm cavity, so the dark-photon field is coherent
across the detector volume --- the haloscope picks up the same
phase everywhere inside the cavity, and the rate is enhanced by a
coherence factor.
\end{solution}
\end{exercise}

\begin{exercise}[Dark-photon absorption in aluminium]
For the dark-photon parameters of \cref{ex:dark-photon-rate}
($m_{A'}=1\,\text{meV}$, $\varepsilon=10^{-13}$, Al absorber), the
absorption cross-section per unit volume is
$\sigma_\text{abs}^\text{ord}\sim 10^4\,\text{cm}^{-1}$ at
$\omega=242\,\text{GHz}$. Verify the rate quoted in the example,
$\sim 10^{-3}\,\text{Hz}/\text{cm}^3$.
\begin{solution}
$\dd N/(\dd t\,\dd V)=\varepsilon^2\rho_\chi\sigma_\text{abs}^\text{ord}/m_{A'}$.
With $\rho_\chi/m_{A'}=0.3\,\text{GeV}/\text{cm}^3 / 10^{-3}\,\text{eV}
=3\times 10^{14}\,\text{cm}^{-3}$, $\varepsilon^2=10^{-26}$, and
$\sigma_\text{abs}^\text{ord}\cdot c=10^4\,\text{cm}^{-1}\cdot 3\times 10^{10}\,\text{cm/s}
=3\times 10^{14}\,\text{cm}^{-1}\,\text{s}^{-1}$ (taking the rate
form $n\,\sigma\,c$ with $n$ the photon-equivalent density), we
obtain $\sim 10^{-26}\cdot 3\times 10^{14}\cdot 3\times 10^{14}\cdot
3\times 10^{-15}\,\text{s}^{-1}\,\text{cm}^{-3}\sim 10^{-3}$,
matching the order-of-magnitude estimate.
\end{solution}
\end{exercise}

\begin{exercise}[Poisson upper limit for low counts]
Solve \cref{eq:poisson-UL} for $\mu_{90}$ at $b=0$ for $N=0,1,2,3$
observed events.
\begin{solution}
For $N=0$: $e^{-\mu_{90}}=0.10\Rightarrow\mu_{90}=2.30$.
For $N=1$: $(1+\mu_{90})e^{-\mu_{90}}=0.10\Rightarrow\mu_{90}\approx 3.89$.
For $N=2$: $(1+\mu_{90}+\mu_{90}^2/2)e^{-\mu_{90}}=0.10\Rightarrow\mu_{90}\approx 5.32$.
For $N=3$: $(1+\mu_{90}+\mu_{90}^2/2+\mu_{90}^3/6)e^{-\mu_{90}}=0.10
\Rightarrow\mu_{90}\approx 6.68$.
\end{solution}
\end{exercise}

\begin{exercise}[Profile-likelihood limit, $b=0$]
Show that for a one-bin Poisson experiment with $b=0$ and $N$
observed, the profile-likelihood test statistic
\cref{eq:wilks} is
$t_\mu=2[\mu-N+N\ln(N/\mu)]$ for $N>0$ and $t_\mu=2\mu$ for $N=0$.
Compute the one-sided 90 \% CL upper limit for $N=0$.
\begin{solution}
The log-likelihood $\ln L=-\mu+N\ln\mu-\ln N!$ has global maximum
at $\hat\mu=N$ for $N>0$ and the boundary maximum at $\hat\mu=0$
for $N=0$. Substituting,
$t_\mu=-2\ln L(\mu)+2\ln L(\hat\mu)=2[\mu-\hat\mu+N\ln(\hat\mu/\mu)]$,
which gives the two cases above. For $N=0$, the one-sided 90\,\%
CL upper limit comes from $t_\mu= 2\mu \leq z_{90}^2 = (1.28)^2$
giving $\mu_{90}\approx 0.82$, smaller than the Poisson UL of
$2.30$ because the test statistic uses the likelihood ratio rather
than a tail-probability cut.
\end{solution}
\end{exercise}

\begin{exercise}[Exposure scaling]
A search reports zero events with $\mathcal E_1=10\,\text{kg-day}$,
giving $\mu_{90}=2.30$ on the count and a corresponding
$\sigma_{n,90}=10^{-39}\,\text{cm}^2$ at fixed $m_\chi$. By how
much does the limit improve at exposure
$\mathcal E_2=100\,\text{kg-day}$, assuming background-free
operation?
\begin{solution}
Background-free: the limit on $\mu$ stays at $2.30$ but the
limit on $\sigma_n$ scales inversely with exposure. Hence
$\sigma_{n,90}\to\sigma_{n,90}\cdot\mathcal E_1/\mathcal E_2=10^{-39}\cdot 10^{-1}
=10^{-40}\,\text{cm}^2$ --- one order of magnitude better at ten
times the exposure. With background, the improvement scales as
$\sqrt{\mathcal E}$ instead of $\mathcal E$, motivating the
``low-background'' frontier.
\end{solution}
\end{exercise}

\begin{exercise}[Designing an axion haloscope]
A QCD axion is to be searched for around $m_a c^2=4\,\mu\text{eV}$
with a $V=1\,\text{L}$ cavity in $B_0=8\,\text{T}$, loaded
$Q_L=10^5$, $C_{010}=0.7$, $\kappa=1/2$, single-photon-counting
efficiency $\eta=0.5$, dark-count rate
$\Gamma_\text{dc}=1\,\text{mHz}$. Compute the minimum
$g_{a\gamma\gamma}$ at which a $5$-day scan would produce
a $\geq 5\sigma$ excess over the dark-count background.
\begin{solution}
The signal photon rate is
$\dot N_\text{sig}=\eta P_\text{sig}/(\hbar\omega_a)$ from
\cref{eq:axion-power}. With $\rho_\chi=0.3\,\text{GeV}/\text{cm}^3$,
$\omega_a/2\pi=1\,\text{GHz}$, plug in the numerical prefactors:
$P_\text{sig}\sim 10^{-22}\,(g_{a\gamma\gamma}/10^{-15}\,\text{GeV}^{-1})^2\,\text{W}$,
giving $\dot N_\text{sig}\sim 10^{-2}\,(g/10^{-15})^2\,\text{photons/s}$.
The background count over $T=5\,\text{days}=4\times 10^5\,\text{s}$
is $N_b=\Gamma_\text{dc}T=400$, with Poisson noise
$\sqrt{N_b}=20$. Demanding signal exceeds $5\sqrt{N_b}=100$ over
$T$:
$\dot N_\text{sig} T\geq 100\Rightarrow (g/10^{-15})^2\geq
2.5\times 10^{-2}/(\text{factor})$, giving roughly
$g_{a\gamma\gamma}\gtrsim 10^{-15}\,\text{GeV}^{-1}$
(neglecting cavity-bandwidth-coverage corrections), at the boundary of the QCD-axion band.
\end{solution}
\end{exercise}

\begin{exercise}[Reading an exclusion curve]
The exclusion curves of \cref{ex:exclusion-sketch} typically run
flat above some characteristic mass and turn upward sharply below
a threshold-driven floor. Explain qualitatively, in terms of
\cref{eq:E_R-max} and \cref{eq:N-exp}, why the curve cuts off at
low mass and how the cutoff scales with the detector threshold
$E_\text{th}$.
\begin{solution}
For $m_\chi\to 0$, $\mu_{\chi T}\to m_\chi$ and
$E_R^\text{max}\propto m_\chi^2$, so the recoil-energy spectrum
shifts to lower and lower energies. Once $E_R^\text{max}<E_\text{th}$
no recoil is detectable, and \cref{eq:N-exp} drops to zero
regardless of cross-section --- there is no useful constraint
on $\sigma_n$ at that mass. The threshold of the exclusion curve
in $m_\chi$ is therefore set by
$m_\chi^\text{min}\sim\sqrt{M_T E_\text{th}}/v_\text{esc}$:
lowering $E_\text{th}$ by a factor of $10$ pushes the reach
downward by $\sqrt{10}\sim 3$. This is the central engineering
incentive of the cQED phonon-detector programme.
\end{solution}
\end{exercise}
